\documentclass[10pt,a4paper]{amsart}
\usepackage[margin=19mm]{geometry}
\usepackage{amsmath,amssymb,mathtools}
\usepackage[scr=rsfs,cal=euler]{mathalfa}
\usepackage{xfrac}
\usepackage[unicode,hidelinks,bookmarks=false]{hyperref}
\newcommand{\ie}{i.e.\ }
\newcommand{\cf}{cf.\ }
\newcommand{\resp}{respectively\ }
\newcommand{\Subsection}{\subsection}
\providecommand{\nobreakdash}{\nobreak\hbox{-}}
\setlength{\emergencystretch}{2em}
\multlinegap0pt
\usepackage{amssymb}
\newtheorem{lemma}{Lemma}
\newtheorem{notation}{Notation}
\newcommand{\End}{{ \rm End }}
\newcommand{\hueca}[1]{\mathbb{#1}}
\newcommand{\rightmap}[1]{\smash{\mathop{\hbox to 
20pt{\rightarrowfill}}\limits^{#1}}}
\title{\bf {\Large On generic bricks over tame algebras}}
\author{R. Bautista, E. P\'erez, L. Salmer\'on}
\date{Source: arXiv:2408.16127v2 (CC BY 4.0)}
\begin{document}
\maketitle

\noindent\emph{Source and licence.} \bf {\Large On generic bricks over tame algebras}. Version arXiv:2408.16127v2 (CC BY 4.0), distributed by arXiv under the Creative Commons Attribution 4.0 International licence, http://creativecommons.org/licenses/by/4.0/, as recorded in the arXiv licence metadata for this exact version and shown on the abstract page https://arxiv.org/abs/2408.16127v2. Full licence notice: \texttt{licenses/arxiv-2408.16127-CC-BY-4.0.txt}.\par\smallskip

\noindent\textit{Editorial scope.} Counted unit: the article's lemma sobre (source0.tex lines 1603--1605). The article's complete original proof of it is delivered with it (source0.tex lines 1607--1680, one \texttt{proof} environment of 74 lines). No mathematics is written, repaired or re-derived: the statement and the proof are the article's own lines, in the article's own notation, and no retained line is substituted or re-flowed.  Retained uncounted as the source-local context the proof needs: lines 1583--1586; lines 1592--1599.  Nothing was omitted from any retained span, so the package carries no omission record.  No retained line contains a citation, so the wrapper carries no bibliography and no bibliography entry.  No retained line contains a figure environment, an image-inclusion command or drawing code, so no image is delivered and the package declares no asset dependency.  Editorial formatting only, in addition: an AMS \texttt{amsart} preamble that restates the article's own declarations and macros used by the retained lines, namely the environments \texttt{lemma}, \texttt{notation} on separate counters, together with the standard packages those lines need.  The retained environments print in this document as Lemma 1, Notation 1 in order of appearance, whereas the article prints the same items with the numbering of its own full document.  The complete native source of the article as distributed in the arXiv e-print for this exact version is bundled unchanged as \texttt{sources/164424/source0.tex}.
\begin{lemma}\label{L: cocientes combin lin de basicos en k(x)} For any $r,s\in \hueca{N}$ and $\lambda\in k$, we have in $k(x)$: 
$$x^r/(x-\lambda)^s=h_{r,s}(x)+u_{r,s}(x)+\lambda^r/(x-\lambda)^s,$$
where $h_{r,s}(x)$ is a polynomial in $k[x]$ and $u_{r,s}(x)\in k(x)$ is a $k$-linear combination of the rational functions $\{1/(x-\lambda)^i\mid 0<i<s\}$. 
\end{lemma}

\begin{notation}\label{N: def de qsubc} Given any $r(x)\in k(x)$, we denote by $\mu_{r(x)}\in \End_k(k(x))$ the linear map given by multiplication by $r(x)$. 
For any linear map $q\in \End_k(k(x))$, we consider the associated morphism of $k$-algebras $k[x,y]\rightmap{}\End_k(k(x))$ such that $x\mapsto q\mu_x$ and $y\mapsto \mu_xq$. 
We denote the image of each $c\in k[x,y]$ by $q_c$. Thus, if $c(x,y)=\sum_{j=0}^ma_j(y)x^j\in k[x,y]$, with $a_j(y)\in k[y]$, and $z\in k(x)$, we have 
$$q_c(z)=\sum_{j=0}^ma_j(x)q(x^jz)=\left[\sum_ja_j(x)qx^j\right](z).$$
Given $c_1,c_2\in k[x,y]$, we have 
$q_{c_1c_2}=q_{c_1}q_{c_2}$. 
 Notice that, if $c,c_1\in k[x,y]$ satisfy $c_1(x,y)=c(x,y)s(x)$, for some non-zero $s(x)\in k[x]$, then the linear endomorphism  $q':=q\mu_{s(x)}^{-1}\in \End_k(k(x))$, is such that   $q'_{c_1}=q_c$. 
\end{notation}

\begin{lemma}[{\bf Convolution Lemma}]\label{L: convolution map sobre} If $0\not=c(x,y)\in k[x,y]$,  the map 
$$\End_k(k(x))\rightmap{}\End_k(k(x))\hbox{ such that }q\mapsto q_c$$ is surjective. 
\end{lemma}

\begin{proof} Assume that $c(x,y)=\sum_{j=0}^ma_j(y)x^j$, with $a_m(y)\not=0$, and take a linear map $p\in \End_k(k(x))$. Let us consider first the following:
\medskip

\noindent\emph{Special Case: There is no $\lambda\in k$ such that 
$c(\lambda,y)=\sum_{j=0}^ma_j(y)\lambda^j=0.$}
\medskip


In order to define the appropriate linear map $q:k(x)\rightmap{}k(x)$, we have to define its values on the $k$-basis of $k(x)$ formed by the elements $x^i$ and $1/(x-\lambda)^j$, for $i\geq 0$, $j\geq 1$, and $\lambda\in k$. We are looking for a linear map $q$ such that the equations 
\begin{equation}\tag{$Eq(z)$}
 p(z)=\sum_{j=0}^ma_j(x)q(x^jz)
\end{equation}
 are satisfied for each one of these basis elements $z$. 



Having in mind the basis element $z=1$, we first define the images $q(1)$, $q(x)$,...,$q(x^m)$ such that the equation $Eq(1)$ is satisfied:   
$$a_0(x)q(1)+a_1(x)q(x)+\cdots+a_m(x)q(x^m)=p(1).$$
For this, we first notice that $c(x,y)\not=0$ implies that $c(x^n,x)\not=0$, for $n\in \hueca{N}$ big enough. Then, we take $q(1):=c(x^n,x)^{-1}p(1)$ and, $q(x^j):=x^{nj}q(1)$, for $j\in [1,m]$. Equation $Eq(1)$ holds because  we have 
$$\begin{matrix}
\sum_{j=0}^ma_j(x)q(x^j)&=&a_0(x)c(x^n,x)^{-1} p(1)+\sum_{j=1}^ma_j(x)x^{nj}c(x^n,x)^{-1}p(1)\hfill\\
&=&
(\sum_{j=0}^ma_j(x)x^{nj})c(x^n,x)^{-1}p(1)=p(1).\hfill\\
\end{matrix}$$

Then, having in mind the basis element $z=x$,  we define 
$$q(x^{m+1}):=a_m(x)^{-1}[p(x)-a_0(x)q(x)-a_1(x)q(x^2)-\cdots -a_{m-1}(x)q(x^m)],$$
so that the equation $Eq(x)$ is satisfied. Similarly, in order to comply with the equations $Eq(x^i)$, we define inductively 
$$q(x^{m+i}):=a_m(x)^{-1}[p(x^i)-a_0(x)q(x^i)-a_1(x)q(x^{i+1})-\cdots -a_{m-1}(x)q(x^{m-1+i})].$$
Next, we want to define the value $q(1/(x-\lambda))$ in order to comply with equation $Eq(1/(x-\lambda))$, that is with 
$$p(1/(x-\lambda))=\sum_{j=0}^ma_j(x)q(x^j/(x-\lambda)).$$
From (\ref{L: cocientes combin lin de basicos en k(x)}), we know that $x^j/(x-\lambda)=h_j(x)+\lambda^j/(x-\lambda)$, for some $h_j(x)\in k[x]$, for each $j\in [1,m]$. So, the preceding equation can be written as 
$$\begin{matrix}
p(1/(x-\lambda))
&=&
\sum_{j=0}^ma_j(x)q[h_j(x)+\lambda^j/(x-\lambda)]\hfill\\
&=&
\sum_{j=0}^ma_j(x)q(h_j(x)) +\sum_{j=0}^ma_j(x)\lambda^jq(1/(x-\lambda))\hfill\\
&=&
\sum_{j=0}^ma_j(x)q(h_j(x)) +c(\lambda,x)q(1/(x-\lambda)).\hfill\\
\end{matrix}
$$
So, if we define 
$$q(1/(x-\lambda)):=c(\lambda,x)^{-1}\left[p(1/(x-\lambda))-\sum_{j=0}^ma_j(x)q(h_j(x))\right],$$
the equation $Eq(1/(x-\lambda))$ is satisfied. Now, assume that $s\geq 1$ and that $q(1/(x-\lambda)^i)$ has been defined for $i\in [1,s]$ such that the corresponding equations $Eq(1/(x-\lambda)^i)$ hold. As in the preceding case, we 
want to define the value $q(1/(x-\lambda)^{s+1})$ in order to comply with equation $Eq(1/(x-\lambda)^{s+1})$, that is with 
$$p(1/(x-\lambda)^{s+1})=\sum_{j=0}^ma_j(x)q(x^j/(x-\lambda)^{s+1}).$$
From (\ref{L: cocientes combin lin de basicos en k(x)}), we know that $x^j/(x-\lambda)^{s+1}=h_j(x)+u_j(x)+\lambda^j/(x-\lambda)^{s+1}$, for some $h_j(x)\in k[x]$ and $u_j(x)\in k(x)$, where $u_j(x)$ is a $k$-linear combination of the rational functions $\{1/(x-\lambda)^i\mid 0<i\leq s\}$, for each $j\in [1,m]$. So, the preceding equation can be written as 
$$\begin{matrix}
p(1/(x-\lambda)^{s+1})
&=&
\sum_{j=0}^ma_j(x)q[h_j(x)+u_j(x)+\lambda^j/(x-\lambda)^{s+1}]\hfill\\
&=&
\sum_{j=0}^ma_j(x)q[h_j(x)+u_j(x)] +c(\lambda,x)q(1/(x-\lambda)^{s+1}).\hfill\\
\end{matrix}
$$
So, if we define 
$$q(1/(x-\lambda)^{s+1}):=c(\lambda,x)^{-1}\left[p(1/(x-\lambda)^{s+1})-\sum_{j=0}^ma_j(x)q[h_j(x)+u_j(x)]\right],$$
the equation $Eq(1/(x-\lambda)^{s+1})$ is satisfied. Thus there is $q\in \End_k(k(x))$ such that $q_c=p$. 
 \medskip

\noindent\emph{General Case:} 
\medskip

By assumption, $0\not=c(x,y)\in k[x,y]$. We claim that there are $\lambda_1,\ldots,\lambda_t\in k$ and a polynomial $c_t(x,y)\in k[x,y]$ such that 
$$c(x,y)=c_t(x,y)(x-\lambda_1)\cdots(x-\lambda_t),$$
and $c_t(\lambda,y)\not=0$, for all $\lambda\in k$. 

Indeed, if $c(\lambda_1,y)=0$, for some $\lambda_1\in k$, the scalar $\lambda_1$ is a root of $c(x,y)\in k[y][x]$, so $c(x,y)=c_1(x,y)(x-\lambda_1)$. If $c_1(\lambda_2,y)=0$, for some $\lambda_2\in k$, we obtain $c(x,y)=c_2(x,y)(x-\lambda_1)(x-\lambda_2)$, with $c_2(x,y)\in k[x,y]$. This process can only be repeated a finite number of times because the degree of the polynomial $c(x,y)\in k[y][x]$ is finite. Our claim follows from this. 


Thus, 
  $c(x,y)=c_t(x,y)s(x)$, with $s(x)=(x-\lambda_1)\cdots(x-\lambda_t)$, for some $\lambda_1,\ldots,\lambda_t\in k$, and $c_t(x,y)\in k[x,y]$ such that $c_t(\lambda,y)\not=0$ for all $\lambda\in k$. Thus, given $p\in \End_k(k(x))$, there is some $q\in \End_k(k(x))$ with $q_{c_t}=p$, by the preceding special case. From the remark at the end of (\ref{N: def de qsubc}), we know that $q':=q\mu^{-1}_{s(x)}\in \End_k(k(x))$ is such that $q'_c=q_{c_t}=p$. 
\end{proof}

\end{document}
